% The DGX Spark Inference Atlas — paper skeleton.
% Headline numbers come from results/derived/macros.tex, which is generated
% EXCLUSIVELY by scripts/sync_readme.py from committed run artifacts. Never
% hand-edit a number in this file; add a macro upstream instead.
\documentclass[sigconf, nonacm]{acmart}

\usepackage{booktabs}
\usepackage{graphicx}

% Generated macros (sync_readme.py). Placeholder keeps the skeleton buildable.
\IfFileExists{../results/derived/macros.tex}{%
  % Auto-generated by scripts/sync_readme.py — do not edit.
\newcommand{\atlasChangepoints}{0}
\newcommand{\atlasCoutCiHiE20}{0.0000}
\newcommand{\atlasCoutCiHiE30}{0.0000}
\newcommand{\atlasCoutCiHiE30A}{0.0000}
\newcommand{\atlasCoutCiHiE30W}{0.0000}
\newcommand{\atlasCoutCiHiE40}{0.0000}
\newcommand{\atlasCoutCiHiFP16}{0.0000}
\newcommand{\atlasCoutCiLoE20}{0.0000}
\newcommand{\atlasCoutCiLoE30}{0.0000}
\newcommand{\atlasCoutCiLoE30A}{0.0000}
\newcommand{\atlasCoutCiLoE30W}{0.0000}
\newcommand{\atlasCoutCiLoE40}{0.0000}
\newcommand{\atlasCoutCiLoFP16}{0.0000}
\newcommand{\atlasCoutE20}{0.0000}
\newcommand{\atlasCoutE30}{0.0000}
\newcommand{\atlasCoutE30A}{0.0000}
\newcommand{\atlasCoutE30W}{0.0000}
\newcommand{\atlasCoutE40}{0.0000}
\newcommand{\atlasCoutFP16}{0.0000}
\newcommand{\atlasMcNemarFamily}{0}
\newcommand{\atlasNE20}{0}
\newcommand{\atlasNE30}{0}
\newcommand{\atlasNE30A}{0}
\newcommand{\atlasNE30W}{0}
\newcommand{\atlasNE40}{0}
\newcommand{\atlasNFP16}{0}
\newcommand{\atlasProvEnvCaptures}{3}
\newcommand{\atlasProvMarketSnapshots}{2}
\newcommand{\atlasProvRunnerRuns}{0}
\newcommand{\atlasProvToolcallRuns}{0}
%
}{%
  % Placeholder macros so the paper builds before any experiment has run.
% The real values are generated into ../results/derived/macros.tex by
% scripts/sync_readme.py and take precedence when present.
\newcommand{\atlasCoutFPXX}{0.0000}
\newcommand{\atlasCoutCiLoFPXX}{0.0000}
\newcommand{\atlasCoutCiHiFPXX}{0.0000}
\newcommand{\atlasCoutE40}{0.0000}
\newcommand{\atlasCoutCiLoE40}{0.0000}
\newcommand{\atlasCoutCiHiE40}{0.0000}
\newcommand{\atlasCoutE30}{0.0000}
\newcommand{\atlasCoutCiLoE30}{0.0000}
\newcommand{\atlasCoutCiHiE30}{0.0000}
\newcommand{\atlasCoutE20}{0.0000}
\newcommand{\atlasCoutCiLoE20}{0.0000}
\newcommand{\atlasCoutCiHiE20}{0.0000}
\newcommand{\atlasMcNemarFamily}{0}
\newcommand{\atlasChangepoints}{0}
%
}

\begin{document}

\title{The DGX Spark Inference Atlas: Quantized Serving of Agentic Workloads
on GB10-Class Hardware}

\author{The Atlas Contributors}
\affiliation{\institution{Open measurement study}}
\email{see repository}

\begin{abstract}
Aggressive low-bit quantization (2--4 bits per weight) is the enabling
condition for serving modern mixture-of-experts models on unified-memory
appliances such as NVIDIA's DGX Spark. Community reports in 2026 disagree
about both quality retention and decode throughput (TensorFold claims
1.6--3$\times$ vLLM; patched-vLLM reports claim the opposite), and existing
benchmarks rarely evaluate \emph{agentic} capability against \emph{live}
external data. This atlas contributes a pre-registered measurement protocol
with (i) live-data-grounded tool-calling evaluation, (ii) isolation-controlled
timing and NVML energy accounting, and (iii) an executable replication kit
for DGX Spark hardware. On the locally executable arm we measure an EXL3
bit-width ladder with \atlasMcNemarFamily{} comparisons against an unquantized
reference; headline C\_outcome at 4\,bpw is \atlasCoutE40{} (95\% CI
\atlasCoutCiLoE40--\atlasCoutCiHiE40) versus \atlasCoutFP16{} for FP16, with
\atlasChangepoints{} detected changepoint(s) on the ladder
($\alpha=0.01$). External Spark claims are catalogued and quarantined
rather than repeated.
\end{abstract}

\keywords{quantization, EXL3, DGX Spark, agentic evaluation, reproducibility}

\maketitle

\section{Introduction}
Deployment of frontier-class open models on single-node appliances has moved
from curiosity to production line in 2026: EXL3 trellis quants of
GLM-5.3-Flash and DeepSeek-v4.1-Flash run on one or two DGX Spars at
13--19\,tok/s decode with speculative drafting. Quality evidence, however,
remains mostly perplexity- and KL-based, while the advertised workloads are
agentic. We ask the questions practitioners actually face: does tool-calling
commitment survive 3\,bpw? Where does it break---in whether-to-call, in
arguments, or in output format? What does a token cost in joules across the
ladder? And can the disputed speeds be adjudicated with a protocol that a
third party can re-run on their own Spark?

\section{Related Work}
We build on benchmarking-methodology guidance \cite{ietf-bench},
sequential-isolation measurement \cite{isolation}, deterministic-first
agentic grading \cite{toolveritas}, McNemar-based degradation detection
\cite{mcnemar}, quantization margin-shrinkage theory \cite{margins,
twomode, divergence, calibrationbottleneck}, energy measurement protocol
\cite{wattcounts}, and the EXL3 \cite{exllamav3} and TensorFold
\cite{tensorfold} serving stacks. Full ecosystem map, including the 2026
Spark speed dispute and its sources, is in the repository's
\texttt{docs/RELATED\_WORK.md}.

\section{Method}
See \texttt{docs/METHODOLOGY.md} (pre-registered 2026-10-18). Conditions:
FP16 reference; EXL3 at 4/3/2\,bpw with shared calibration; a
calibration-corpus ablation at 3\,bpw. Protocol: IETF warm-up; sequential
isolation; fixed-shape prompts; streaming TTFT; per-condition NVML energy;
live Kraken/Frankfurter tool execution during evaluation with episode-time
truth capture; deterministic-first three-cell grading with repair flags.
Analysis: McNemar exact vs reference, bootstrap CIs, TOST for equivalence,
permutation changepoints on the bpw axis, Holm--Bonferroni across the
family.

\section{Results}
% All numbers below expand from generated macros; empty macros compile as 0.
Table placeholders follow; the committed markdown tables in
\texttt{results/derived/tables/} are the canonical rendering.

\begin{table}[h]
\caption{Tool-calling outcomes on live external data (C\_outcome).}
\begin{tabular}{lrr}
\toprule
Condition & C\_outcome & 95\% CI \\
\midrule
FP16 & \atlasCoutFP16 & \atlasCoutCiLoFP16--\atlasCoutCiHiFP16 \\
EXL3 4\,bpw & \atlasCoutE40 & \atlasCoutCiLoE40--\atlasCoutCiHiE40 \\
EXL3 3\,bpw & \atlasCoutE30 & \atlasCoutCiLoE30--\atlasCoutCiHiE30 \\
EXL3 2\,bpw & \atlasCoutE20 & \atlasCoutCiLoE20--\atlasCoutCiHiE20 \\
\bottomrule
\end{tabular}
\end{table}

Energy per token by condition is reported in
\texttt{results/derived/tables/timing.md}; conditions flagged by the
isolation gate (CV $>0.10$) are marked and excluded from ladder fitting.

\section{Limitations}
Local-arm model scale is small (1.5B-class); absolute quality numbers do not
transfer to 355B-class MoE models, the \emph{method} does. N per condition
bounds detectable degradation (power analysis in METHODOLOGY \S6);
non-significance is never reported as equivalence without TOST. NVML energy
is whole-device and protocol-relative. DGX Spark (GB10) numbers from
community sources are labelled external and excluded from measured tables;
the Spark Kit (\texttt{docs/SPARK\_KIT.md}) exists to convert them into
measured ones.

\section{Reproducibility}
The repository ships a docker-compose reproducibility floor (CPU-only smoke
path), pinned harness digests in every artifact, and a manifest
(\texttt{results/derived/readme-sync.json}) binding this paper's macros to
the exact raw records. See \texttt{papers/ARTIFACT\_APPENDIX.md} following
PACT-2026-AE conventions.

\bibliographystyle{ACM-Reference-Format}
\bibliography{references}

\end{document}
